% Gerado por scripts/migrate_markdown_to_latex.py.
% Fonte migrada: fases/01_validacao_conceitual/docs/base_teorica_validacao.md


\chapter{Base teorica e validacao cientifica do modelo}



\section{Objetivo do documento}





Este documento organiza a primeira fase do trabalho: estabelecer uma base teorica confiavel e um metodo de validacao antes de qualquer benchmark de hardware. A regra de trabalho sera:



\begin{icquote}
Nenhum conceito tecnico deve ser usado no software ou no relatorio sem uma referencia cientifica confiavel, uma definicao operacional e um teste que verifique se a implementacao observada corresponde ao comportamento esperado.
\end{icquote}







Isso evita um erro comum em experimentos de software: medir desempenho de uma implementacao que usa o nome de uma tecnica, mas nao realiza a tecnica de forma cientificamente adequada. Por exemplo, uma funcao que quantiza e depois volta imediatamente para \texttt{float32} pode ser util para estudar erro numerico, mas nao deve ser apresentada como evidencia de ganho real de hardware.



\section{Principio metodologico}



O trabalho deve separar quatro niveis de certeza:


\begin{enumerate}
\item Definicao: o termo usado esta definido corretamente pela literatura?
\item Matematica: a formula implementada bate com casos pequenos derivados da
definicao?
\item Algoritmo: o software implementa o algoritmo, ou uma aproximacao aceitavel e
explicitamente justificada?
\item Hardware: a implementacao realmente usa uma representacao,
estrutura de dados, kernel ou caminho de execucao capaz de produzir ganho
fisico mensuravel em memoria, latencia ou throughput?
\end{enumerate}





Uma implementacao pode ser valida em um nivel e invalida em outro. Por exemplo, zerar 50\% dos pesos de uma matriz e continuar executando uma multiplicacao densa pode ser um pruning conceitual/algoritmico, mas nao prova ganho real de hardware.



\section{Matriz de rastreabilidade}




Todo conceito, ferramenta ou algoritmo usado no projeto deve entrar em uma matriz com os seguintes campos.



\begin{longtable}{@{}>{\RaggedRight\arraybackslash}p{0.410\linewidth}>{\RaggedRight\arraybackslash}p{0.410\linewidth}@{}}
\toprule
Campo & Descricao \\
\midrule
\endfirsthead
\toprule
Campo & Descricao \\
\midrule
\endhead
Conceito declarado & Nome usado no trabalho, como Transformer, quantizacao ou pruning. \\
Fonte cientifica & Artigo, livro-texto ou documentacao primaria que define o conceito. \\
Definicao aceita & Definicao curta usada pelo projeto. \\
Propriedade esperada no modelo & O que precisa aparecer na implementacao para o nome ser valido. \\
Teste de verificacao & Como verificar o comportamento esperado. \\
Resultado observado & Campo a preencher apos executar a verificacao. \\
Status & Aderente, parcialmente aderente ou nao aderente. \\
Nivel de validade & Conceitual, numerico, algoritmico ou hardware. \\
\bottomrule
\end{longtable}



Modelo de linha:



\begin{landscape}
\scriptsize
\begin{longtable}{@{}>{\RaggedRight\arraybackslash}p{0.102\linewidth}>{\RaggedRight\arraybackslash}p{0.102\linewidth}>{\RaggedRight\arraybackslash}p{0.102\linewidth}>{\RaggedRight\arraybackslash}p{0.102\linewidth}>{\RaggedRight\arraybackslash}p{0.102\linewidth}>{\RaggedRight\arraybackslash}p{0.102\linewidth}>{\RaggedRight\arraybackslash}p{0.102\linewidth}>{\RaggedRight\arraybackslash}p{0.102\linewidth}@{}}
\toprule
Conceito declarado & Fonte cientifica & Definicao aceita & Propriedade esperada no modelo & Teste de verificacao & Resultado observado & Status & Nivel \\
\midrule
\endfirsthead
\toprule
Conceito declarado & Fonte cientifica & Definicao aceita & Propriedade esperada no modelo & Teste de verificacao & Resultado observado & Status & Nivel \\
\midrule
\endhead
Transformer & Vaswani et al. (2017) & Arquitetura baseada em atencao, com tratamento de sequencias e informacao posicional. & Entrada sequencial, Q/K/V, atencao escalonada, dependencia entre posicoes e informacao posicional ou justificativa formal de ausencia. & Comparar a saida de um caso pequeno com implementacao de referencia e auditar componentes arquiteturais. & A preencher. & A preencher. & Conceitual/\allowbreak{}algoritmico. \\
\bottomrule
\end{longtable}
\end{landscape}



\section{Glossario operacional}



\subsection{Rede neural}


\begin{itemize}
\item Definicao cientifica: modelo computacional parametrizado composto por
unidades ou camadas que aplicam transformacoes lineares e nao lineares para
aproximar funcoes a partir de dados.
\item Referencia confiavel: Goodfellow, Bengio e Courville, \emph{Deep Learning},
especialmente a discussao de redes feedforward e modelos parametrizados:
.
\item Explicacao intuitiva: uma rede neural e uma funcao com muitos parametros
ajustaveis. Ela transforma uma entrada em uma saida por etapas.
\item Uso no trabalho: sera o conceito mais geral por tras dos modelos estudados.
Um Transformer e um tipo especifico de arquitetura neural.
\item Propriedades exigidas no codigo: parametros numericos, fluxo de entrada para
saida, composicao de operacoes e possibilidade de inferencia com parametros
fixos.
\item Teste de aderencia: verificar se a implementacao possui parametros, executa
uma passagem direta deterministica para uma entrada fixa e produz saida com
dimensao esperada.
\end{itemize}


\subsection{Deep learning}


\begin{itemize}
\item Definicao cientifica: subarea de aprendizado de maquina baseada em multiplas
camadas de representacao, nas quais conceitos mais complexos sao construidos a
partir de conceitos mais simples.
\item Referencia confiavel: Goodfellow, Bengio e Courville, \emph{Deep Learning}:
.
\item Explicacao intuitiva: deep learning e o uso de redes neurais com varias
etapas de transformacao, permitindo aprender representacoes intermediarias.
\item Uso no trabalho: delimita o contexto em que Transformers, compressao e
inferencia sao discutidos.
\item Propriedades exigidas no codigo: mais de uma transformacao parametrizada ou
uma arquitetura explicitamente derivada de modelos profundos.
\item Teste de aderencia: inspecionar a arquitetura. Se houver apenas uma
multiplicacao linear isolada, o objeto deve ser descrito como uma operacao
auxiliar, nao como modelo de deep learning completo.
\end{itemize}


\subsection{Transformer}


\begin{itemize}
\item Definicao cientifica: arquitetura neural baseada em mecanismos de atencao,
proposta por Vaswani et al. (2017), dispensando recorrencia e convolucoes na
arquitetura original de sequencia para sequencia.
\item Referencia confiavel: Vaswani et al., \emph{Attention Is All You Need}, 2017:
.
\item Explicacao intuitiva: um Transformer processa uma sequencia permitindo que
cada posicao consulte outras posicoes por meio de atencao.
\item Uso no trabalho: sera o objeto arquitetural de referencia. O projeto podera
estudar uma operacao inspirada em Transformer, um bloco Transformer
simplificado ou uma implementacao Transformer completa.
\item Propriedades exigidas no codigo: entrada sequencial, projecoes Q/K/V,
mecanismo de atencao, mistura de informacao entre posicoes, informacao
posicional ou justificativa formal para remove-la, e transformacoes
parametrizadas.
\item Teste de aderencia: executar checklist arquitetural e comparar um caso
pequeno com uma implementacao de referencia. Se faltar componente essencial,
classificar corretamente como 'operacao inspirada em Transformer' ou 'bloco
Transformer simplificado'.
\end{itemize}


\subsection{Self-attention}


\begin{itemize}
\item Definicao cientifica: mecanismo de atencao no qual consultas, chaves e valores
sao derivados da mesma sequencia, permitindo que cada posicao combine
informacao de outras posicoes da propria entrada.
\item Referencia confiavel: Vaswani et al. (2017):
.
\item Explicacao intuitiva: cada token 'olha' para os outros tokens da mesma
sequencia e decide quais importam mais para sua representacao.
\item Uso no trabalho: sera uma das operacoes centrais a validar antes de estudar
custo computacional.
\item Propriedades exigidas no codigo: Q, K e V calculados da mesma entrada, matriz
de atencao com dimensao compativel com o comprimento da sequencia e saida
dependente de multiplas posicoes.
\item Teste de aderencia: alterar um token nao trivial da sequencia e verificar se
a saida de outras posicoes pode mudar. Comparar tambem a saida com uma
implementacao direta da formula de atencao.
\end{itemize}


\subsection{Scaled dot-product attention}


\begin{itemize}
\item Definicao cientifica: atencao definida por
\texttt{Attention(Q, K, V) = softmax(QK\textasciicircum{}T /\allowbreak{} sqrt(d\_\allowbreak{}k))V}.
\item Referencia confiavel: Vaswani et al. (2017):
.
\item Explicacao intuitiva: calcula similaridades entre consultas e chaves,
normaliza essas similaridades e usa os pesos resultantes para combinar os
valores.
\item Uso no trabalho: sera a forma padrao para validar a operacao de atencao.
\item Propriedades exigidas no codigo: produto \texttt{QK\textasciicircum{}T}, divisao por \texttt{sqrt(d\_\allowbreak{}k)},
\texttt{softmax} no eixo correto e multiplicacao pelos valores \texttt{V}.
\item Teste de aderencia: criar tensores pequenos com valores conhecidos, calcular
a saida manualmente ou por uma implementacao de referencia e exigir diferenca
numerica dentro de uma tolerancia definida.
\end{itemize}


\subsection{Inferencia}


\begin{itemize}
\item Definicao cientifica: fase em que um modelo ja definido ou treinado e usado
para produzir saidas a partir de entradas, sem atualizar parametros.
\item Referencia confiavel: Goodfellow, Bengio e Courville, \emph{Deep Learning}:
. Para uso pratico em software, a
documentacao do PyTorch tambem distingue execucao de modelos em modo de
avaliacao/inferencia.
\item Explicacao intuitiva: inferencia e usar o modelo, nao treina-lo.
\item Uso no trabalho: o foco sera medir custo e comportamento durante execucao de
modelos ou operacoes com pesos fixos.
\item Propriedades exigidas no codigo: ausencia de atualizacao de pesos, modo de
avaliacao quando aplicavel, desativacao de gradientes quando o objetivo for
apenas execucao e entradas controladas.
\item Teste de aderencia: executar a mesma entrada duas vezes com os mesmos pesos e
semente. Os pesos nao devem mudar e a saida deve ser reprodutivel, salvo
operacoes explicitamente nao deterministicas.
\end{itemize}


\subsection{Projecao densa}


\begin{itemize}
\item Definicao cientifica: transformacao linear parametrizada do tipo
\texttt{Y = XW\textasciicircum{}T + b}, geralmente aplicada por multiplicacao matricial densa.
\item Referencia confiavel: Goodfellow, Bengio e Courville para camadas lineares em
redes neurais; Vaswani et al. (2017) para o uso de projecoes aprendidas em Q,
K, V e saidas de atencao.
\item Explicacao intuitiva: uma matriz de pesos transforma cada vetor de entrada em
outro vetor.
\item Uso no trabalho: sera um bloco basico para construir Q, K, V e camadas
lineares simplificadas.
\item Propriedades exigidas no codigo: matriz de pesos densa, dimensoes
consistentes e multiplicacao matricial equivalente a transformacao linear.
\item Teste de aderencia: comparar a saida da funcao com uma multiplicacao
matricial direta em tensores pequenos; verificar dimensoes de entrada, pesos e
saida.
\end{itemize}


\subsection{KV cache}


\begin{itemize}
\item Definicao cientifica: armazenamento de chaves e valores ja calculados durante
geracao autoregressiva, evitando recomputacao em passos futuros.
\item Referencia confiavel: Google Research, *TurboQuant: Redefining AI efficiency
with extreme compression*, 2026:
.
O paper TurboQuant tambem trata quantizacao de vetores e KV cache:
.
\item Explicacao intuitiva: o modelo guarda informacao intermediaria de tokens
anteriores para nao recalcular tudo a cada novo token.
\item Uso no trabalho: sera relevante quando o experimento sair de atencao em lote
para cenarios autoregressivos ou compressao de memoria.
\item Propriedades exigidas no codigo: armazenamento persistente de K e V entre
passos de geracao, crescimento com o comprimento do contexto e reutilizacao em
novos passos.
\item Teste de aderencia: comparar saida com e sem cache para a mesma sequencia. As
saidas devem ser numericamente equivalentes dentro de tolerancia, mas a versao
com cache deve evitar recomputacao de K/V anteriores.
\end{itemize}


\subsection{Pruning}


\begin{itemize}
\item Definicao cientifica: tecnica de compressao que remove parametros, conexoes,
canais ou estruturas consideradas menos relevantes, produzindo esparsidade ou
reducao estrutural.
\item Referencia confiavel: Tang et al., \emph{A Survey on Transformer Compression},
2024: . Para implementacao em ferramenta, a
documentacao oficial do PyTorch descreve \texttt{torch.\allowbreak{}nn.\allowbreak{}utils.\allowbreak{}prune}:
.
\item Explicacao intuitiva: pruning corta partes do modelo que parecem pouco
importantes.
\item Uso no trabalho: sera uma tecnica candidata para comparacao com quantizacao.
\item Propriedades exigidas no codigo: criterio de importancia, mascara ou remocao
estrutural, taxa de poda mensuravel e impacto verificavel na saida.
\item Teste de aderencia: medir fracao de pesos/estruturas removidas, comparar
saida com baseline e verificar se a execucao usa representacao esparsa ou
estrutura reduzida. Se apenas zerar pesos em matriz densa, classificar como
pruning conceitual, nao como ganho de hardware.
\end{itemize}


\subsection{Sparsity}


\begin{itemize}
\item Definicao cientifica: grau de esparsidade de uma representacao, normalmente
associado a proporcao de elementos nulos ou conexoes ausentes em tensores,
matrizes ou estruturas do modelo.
\item Referencia confiavel: Tang et al. (2024) no contexto de compressao de
Transformers; tutorial oficial de pruning do PyTorch para esparsificacao de
redes neurais.
\item Explicacao intuitiva: uma estrutura esparsa tem muitos zeros ou muitas
conexoes removidas.
\item Uso no trabalho: sera uma propriedade esperada de modelos podados.
\item Propriedades exigidas no codigo: contagem objetiva de zeros, mascaras ou
estruturas removidas; quando houver alegacao de hardware, uso de formato ou
kernel que explore esparsidade.
\item Teste de aderencia: calcular \texttt{zeros /\allowbreak{} total} e auditar se a multiplicacao
realmente deixou de processar os elementos removidos. Se o custo computacional
permanecer denso, a validade sera apenas conceitual ou numerica.
\end{itemize}


\subsection{Quantizacao}


\begin{itemize}
\item Definicao cientifica: tecnica de representacao que mapeia valores de maior
precisao para um conjunto menor de niveis discretos, reduzindo bits por valor
e introduzindo erro controlado.
\item Referencia confiavel: Tang et al. (2024) para quantizacao em compressao de
Transformers; Zandieh et al., \emph{TurboQuant}, 2025:
; documentacao \texttt{torchao} para fluxos e
tipos de quantizacao em PyTorch:
.
\item Explicacao intuitiva: quantizar e representar numeros com menos precisao para
economizar memoria e potencialmente acelerar computacao.
\item Uso no trabalho: sera uma tecnica candidata para reduzir custo de inferencia.
\item Propriedades exigidas no codigo: escala ou codigo de quantizacao, valores em
menor precisao ou tipo derivado, erro de dequantizacao mensuravel e, para
hardware, kernel ou representacao que reduza custo real.
\item Teste de aderencia: verificar dtype/bit-width, memoria ocupada, erro em
relacao ao baseline e caminho de execucao. Se os tensores forem
imediatamente convertidos para \texttt{float32}, classificar como simulacao numerica
de quantizacao.
\end{itemize}


\subsection{Latencia}


\begin{itemize}
\item Definicao cientifica: tempo decorrido para completar uma operacao, requisicao
ou unidade de trabalho.
\item Referencia confiavel: Hennessy e Patterson, *Computer Architecture: A
Quantitative Approach*, texto classico de arquitetura de computadores.
\item Explicacao intuitiva: latencia e quanto tempo uma execucao demora.
\item Uso no trabalho: sera uma metrica de desempenho por operacao ou por passo de
inferencia.
\item Propriedades exigidas no codigo: medicao com relogio adequado, aquecimento,
repeticoes, sincronizacao de GPU quando aplicavel e reporte de media e
percentis.
\item Teste de aderencia: executar benchmark com entradas fixas, descartar
aquecimento, sincronizar antes/depois da medicao em GPU e reportar media, p50
e p95.
\end{itemize}


\subsection{Throughput}


\begin{itemize}
\item Definicao cientifica: quantidade de trabalho completada por unidade de tempo,
tambem associada a bandwidth em sistemas computacionais.
\item Referencia confiavel: Hennessy e Patterson, *Computer Architecture: A
Quantitative Approach*.
\item Explicacao intuitiva: throughput mede quanto o sistema processa por segundo.
\item Uso no trabalho: podera ser medido como tokens/s, amostras/s ou operacoes/s.
\item Propriedades exigidas no codigo: definicao explicita da unidade de trabalho e
medicao em janela suficientemente estavel.
\item Teste de aderencia: calcular \texttt{unidades\_\allowbreak{}processadas /\allowbreak{} tempo\_\allowbreak{}total} para lote,
sequencia ou tokens definidos. Nao misturar throughput com latencia sem
indicar a unidade.
\end{itemize}


\subsection{Memoria}


\begin{itemize}
\item Definicao cientifica: recurso de armazenamento usado por parametros,
ativacoes, tensores intermediarios, caches e estruturas auxiliares.
\item Referencia confiavel: Hennessy e Patterson para memoria em sistemas
computacionais; Tang et al. (2024) para reducao de memoria por compressao de
Transformers.
\item Explicacao intuitiva: memoria e o espaco necessario para guardar modelo,
entradas e resultados intermediarios.
\item Uso no trabalho: sera uma metrica central para comparar baseline, pruning e
quantizacao.
\item Propriedades exigidas no codigo: estimativa teorica de bytes e medicao de
pico real quando possivel.
\item Teste de aderencia: calcular \texttt{num\_\allowbreak{}elementos * bytes\_\allowbreak{}por\_\allowbreak{}elemento} para
tensores relevantes e comparar com medidores de alocacao do framework. Para
KV cache, verificar crescimento com o comprimento do contexto.
\end{itemize}


\subsection{FLOPs}


\begin{itemize}
\item Definicao cientifica: contagem de operacoes aritmeticas de ponto flutuante,
usada como estimativa analitica de custo computacional.
\item Referencia confiavel: Hennessy e Patterson para avaliacao quantitativa de
desempenho; Vaswani et al. (2017) para analise de complexidade de atencao e
arquiteturas sequenciais.
\item Explicacao intuitiva: FLOPs estimam quantas contas numericas o algoritmo faz.
\item Uso no trabalho: servira como custo teorico, separado da medicao real de
tempo.
\item Propriedades exigidas no codigo: formula documentada por operacao e dimensoes
usadas no calculo.
\item Teste de aderencia: derivar FLOPs esperados para multiplicacoes matriciais e
atencao; validar com casos pequenos e manter claro que FLOPs nao equivalem
automaticamente a latencia real.
\end{itemize}


\subsection{MSE}


\begin{itemize}
\item Definicao cientifica: erro quadratico medio, calculado pela media dos
quadrados das diferencas entre valores de referencia e valores estimados.
\item Referencia confiavel: Hastie, Tibshirani e Friedman, *The Elements of
Statistical Learning*: . Para
definicao operacional em software, documentacao \texttt{mean\_\allowbreak{}squared\_\allowbreak{}error} do
scikit-learn:
.
\item Explicacao intuitiva: mede erro medio dando peso maior a erros grandes.
\item Uso no trabalho: comparara saidas otimizadas contra a saida baseline.
\item Propriedades exigidas no codigo: comparacao ponto a ponto entre saida de
referencia e saida candidata, com mesmo shape.
\item Teste de aderencia: MSE deve ser zero para vetores identicos e positivo para
diferencas. Testar com exemplos pequenos conhecidos.
\end{itemize}


\subsection{MAE}


\begin{itemize}
\item Definicao cientifica: erro absoluto medio, calculado pela media dos valores
absolutos das diferencas entre referencia e estimativa.
\item Referencia confiavel: Hastie, Tibshirani e Friedman, *The Elements of
Statistical Learning*; definicao operacional \texttt{mean\_\allowbreak{}absolute\_\allowbreak{}error} do
scikit-learn:
.
\item Explicacao intuitiva: mede, em media, o tamanho do erro sem elevar ao
quadrado.
\item Uso no trabalho: complementara o MSE, por ser menos sensivel a erros extremos.
\item Propriedades exigidas no codigo: comparacao ponto a ponto com mesmo shape.
\item Teste de aderencia: MAE deve ser zero para vetores identicos e deve bater com
calculo manual em exemplos pequenos.
\end{itemize}


\subsection{R2}


\begin{itemize}
\item Definicao cientifica: coeficiente de determinacao, usado para avaliar quanto
da variacao da referencia e explicada pela estimativa em relacao a um baseline
medio.
\item Referencia confiavel: Hastie, Tibshirani e Friedman; definicao operacional
\texttt{r2\_\allowbreak{}score} do scikit-learn:
.
\item Explicacao intuitiva: indica se a saida candidata acompanha bem a variacao da
saida de referencia.
\item Uso no trabalho: medira preservacao global do comportamento da saida em
relacao ao baseline.
\item Propriedades exigidas no codigo: referencia e candidato com mesma dimensao;
tratamento explicito de casos constantes.
\item Teste de aderencia: R2 deve ser 1 para previsoes perfeitas, 0 para predicao
igual a media em casos nao constantes e pode ser negativo quando o candidato e
pior que esse baseline.
\end{itemize}


\subsection{Similaridade de cosseno}


\begin{itemize}
\item Definicao cientifica: similaridade entre vetores calculada como produto
interno normalizado pelas normas dos vetores.
\item Referencia confiavel: Manning, Raghavan e Schutze, *Introduction to
Information Retrieval*, Cambridge University Press, 2008:
. Definicao
operacional \texttt{cosine\_\allowbreak{}similarity} do scikit-learn:
.
\item Explicacao intuitiva: mede se dois vetores apontam em direcao parecida,
independentemente da escala.
\item Uso no trabalho: avaliara se a direcao da representacao foi preservada apos
pruning ou quantizacao.
\item Propriedades exigidas no codigo: achatamento ou agregacao definida dos
vetores, tratamento de vetores nulos e formula normalizada.
\item Teste de aderencia: vetores identicos nao nulos devem ter similaridade 1;
vetores ortogonais devem ter similaridade 0; vetores opostos devem tender a
-1.
\end{itemize}


\section{Criterios para chamar uma implementacao de Transformer}




Uma implementacao so deve ser chamada de Transformer se atender aos criterios minimos abaixo ou justificar formalmente qualquer variacao.



\begin{longtable}{@{}>{\RaggedRight\arraybackslash}p{0.410\linewidth}>{\RaggedRight\arraybackslash}p{0.410\linewidth}@{}}
\toprule
Criterio & Verificacao minima \\
\midrule
\endfirsthead
\toprule
Criterio & Verificacao minima \\
\midrule
\endhead
Entrada sequencial & Recebe sequencia de vetores ou embeddings com dimensoes explicitas. \\
Q/K/V & Projeta a entrada em consultas, chaves e valores. \\
Atencao escalonada & Implementa \texttt{softmax(QK\textasciicircum{}T /\allowbreak{} sqrt(d\_\allowbreak{}k))V} ou documenta uma alternativa. \\
Dependencia entre posicoes & A saida de uma posicao pode depender de outras posicoes da sequencia. \\
Informacao posicional & Usa encoding/embedding posicional ou justifica por que o recorte nao precisa dele. \\
Componentes do bloco & Inclui ou justifica ausencia de residual, normalizacao e feed-forward quando o nome usado for bloco Transformer. \\
Referencia numerica & Tem saida comparada contra implementacao direta, biblioteca confiavel ou caso derivado manualmente. \\
\bottomrule
\end{longtable}



Classificacao correta:


\begin{itemize}
\item Operacao inspirada em Transformer: quando apenas uma operacao isolada for
testada, como projecao densa ou scaled dot-product attention.
\item Bloco Transformer simplificado: quando houver atencao, projecoes e parte dos
componentes arquiteturais, mas com simplificacoes assumidas.
\item Transformer: apenas quando a arquitetura tiver os elementos essenciais
definidos pela literatura ou uma justificativa formal para cada diferenca.
\end{itemize}


\section{Validacao de ferramentas e algoritmos}




Antes de usar uma ferramenta em benchmark, ela deve ser validada por tres perguntas:


\begin{enumerate}
\item O que a ferramenta afirma implementar?
\item Qual referencia cientifica ou documentacao primaria sustenta essa afirmacao?
\item O comportamento observado confirma a afirmacao em nivel conceitual,
numerico, algoritmico ou de hardware?
\end{enumerate}


Matriz inicial de ferramentas candidatas:



\begin{landscape}
\scriptsize
\begin{longtable}{@{}>{\RaggedRight\arraybackslash}p{0.164\linewidth}>{\RaggedRight\arraybackslash}p{0.164\linewidth}>{\RaggedRight\arraybackslash}p{0.164\linewidth}>{\RaggedRight\arraybackslash}p{0.164\linewidth}>{\RaggedRight\arraybackslash}p{0.164\linewidth}@{}}
\toprule
Ferramenta/algoritmo & Conceito declarado & Fonte de referencia & Verificacao necessaria & Classificacao esperada \\
\midrule
\endfirsthead
\toprule
Ferramenta/algoritmo & Conceito declarado & Fonte de referencia & Verificacao necessaria & Classificacao esperada \\
\midrule
\endhead
Implementacao manual de atencao & Scaled dot-product attention & Vaswani et al. (2017) & Comparar com calculo manual e/ou biblioteca de referencia. & Conceitual/algoritmica. \\
\texttt{torch.\allowbreak{}nn.\allowbreak{}MultiheadAttention} ou equivalente & Multi-head attention & Documentacao PyTorch + Vaswani et al. & Verificar shapes, mascaras, projecoes e saida contra referencia. & Algoritmica. \\
\texttt{torch.\allowbreak{}nn.\allowbreak{}utils.\allowbreak{}prune} & Pruning/esparsificacao & Tutorial oficial PyTorch + Tang et al. & Medir esparsidade, mascara e impacto na saida; verificar se ha ganho real de execucao. & Conceitual/algoritmica; hardware somente se houver kernel/formato adequado. \\
\texttt{torchao} & Quantizacao & Documentacao \texttt{torchao} + Tang et al. & Verificar dtype, bit-width, escalas, memoria e kernel usado. & Numerica/hardware se usar caminho eficiente. \\
TurboQuant & Quantizacao vetorial/KV cache & Zandieh et al. (2025), Google Research (2026) & Confirmar se ha implementacao disponivel e se o metodo usado corresponde ao paper. & Algoritmica/hardware apenas apos validacao. \\
\bottomrule
\end{longtable}
\end{landscape}



Regras de classificacao:


\begin{itemize}
\item Quantizacao valida para hardware exige representacao menor, kernel compativel
ou reducao real de memoria. Caso contrario, e simulacao numerica.
\item Pruning valido para hardware exige que a esparsidade seja explorada pela
execucao. Caso contrario, e pruning conceitual.
\item Uma biblioteca so sera aceita quando sua documentacao ou artigo explicar
claramente o metodo usado e quando o comportamento observado bater com a
definicao esperada.
\end{itemize}


\section{Procedimento inicial de trabalho}


\begin{enumerate}
\item Revisar e completar este glossario com citacoes formatadas no padrao exigido
pelo relatorio.
\item Escolher uma implementacao de referencia para atencao/Transformer.
\item Criar testes pequenos e deterministas para Q/K/V, atencao e metricas.
\item Validar ferramentas candidatas de pruning e quantizacao contra a matriz.
\item Somente depois disso, definir benchmarks de hardware.
\end{enumerate}



O protocolo detalhado de passagem entre validacao e benchmark esta em \texttt{fases/\allowbreak{}01\_\allowbreak{}validacao\_\allowbreak{}conceitual/\allowbreak{}docs/\allowbreak{}protocolo\_\allowbreak{}experimental.\allowbreak{}md}.




As perguntas de pesquisa, hipoteses e criterios de conclusao estao em \texttt{fases/\allowbreak{}01\_\allowbreak{}validacao\_\allowbreak{}conceitual/\allowbreak{}docs/\allowbreak{}perguntas\_\allowbreak{}pesquisa.\allowbreak{}md}.




A validacao especifica da rede neural usada como bloco Transformer esta em \texttt{fases/\allowbreak{}01\_\allowbreak{}validacao\_\allowbreak{}conceitual/\allowbreak{}docs/\allowbreak{}validacao\_\allowbreak{}transformer.\allowbreak{}md}.



\section{Criterios de aceite desta fase}


\begin{itemize}
\item Todo termo tecnico usado no trabalho possui referencia confiavel.
\item Toda definicao possui uma propriedade verificavel no modelo ou software.
\item Toda ferramenta futura pode ser avaliada pela matriz conceito-fonte-teste.
\item O trabalho consegue dizer, com evidencia, se uma implementacao e valida para
estudo conceitual, teste numerico, teste algoritmico ou teste real de
hardware.
\item Nenhuma conclusao sobre desempenho e apresentada antes da validacao
conceitual e algoritmica.
\end{itemize}


\section{Referencias iniciais}


\begin{itemize}
\item Goodfellow, I.; Bengio, Y.; Courville, A. \emph{Deep Learning}. MIT Press, 2016.
Disponivel em: .
\item Vaswani, A. et al. \emph{Attention Is All You Need}. arXiv:1706.03762, 2017.
Disponivel em: .
\item Tang, Y. et al. \emph{A Survey on Transformer Compression}. arXiv:2402.05964,
\begin{enumerate}
\item Disponivel em: .
\end{enumerate}
\item Zandieh, A. et al. *TurboQuant: Online Vector Quantization with Near-optimal
Distortion Rate*. arXiv:2504.19874, 2025. Disponivel em:
.
\item Google Research. *TurboQuant: Redefining AI efficiency with extreme
compression*. 2026. Disponivel em:
.
\item PyTorch. \emph{Pruning Tutorial}. Disponivel em:
.
\item PyTorch. \emph{Quantization Overview - torchao}. Disponivel em:
.
\item Hastie, T.; Tibshirani, R.; Friedman, J. *The Elements of Statistical
Learning*. Springer, 2009. Disponivel em:
.
\item Manning, C. D.; Raghavan, P.; Schutze, H. *Introduction to Information
Retrieval*. Cambridge University Press, 2008. Disponivel em:
.
\item Hennessy, J. L.; Patterson, D. A. *Computer Architecture: A Quantitative
Approach*. Morgan Kaufmann. Referencia base para latencia, throughput,
memoria e avaliacao quantitativa de desempenho.
\item scikit-learn. \emph{Model evaluation and metrics}. Disponivel em:
.
\end{itemize}
